\documentclass[11pt,a4paper]{article}

\usepackage[a4paper,margin=1in]{geometry}
\usepackage{amsmath}
\usepackage{booktabs}
\usepackage{array}
\usepackage{enumitem}
\usepackage{fancyhdr}
\usepackage[T1]{fontenc}
\usepackage{lmodern}

\newcommand{\universityname}{UNIVERSITY OF DERBY}
\newcommand{\college}{College of Science and Engineering}
\newcommand{\examtitle}{Practice Paper B Solutions}
\newcommand{\moduletitle}{OPERATIONAL RESEARCH}
\newcommand{\modulecode}{5CM522}

\pagestyle{fancy}
\fancyhf{}
\fancyfoot[C]{\thepage}
\renewcommand{\headrulewidth}{0pt}

\setlength{\parindent}{0pt}
\setlength{\parskip}{0.45em}
\setlist[enumerate,1]{label=(\alph*), leftmargin=1.1cm, itemsep=0.4em, topsep=0.3em}

\begin{document}

\begin{center}
{\large \textbf{\universityname}}\\
{\college}\\[0.4em]
{\Large \textbf{\examtitle}}\\
{\large \textbf{\moduletitle} -- \textbf{\modulecode}}
\end{center}

\textbf{1. Scheduling}

\begin{enumerate}
\item Applying Johnson's Rule gives the sequence
\[
J_4,\ J_2,\ J_5,\ J_1,\ J_3
\]

Main steps:

\begin{center}
\begin{tabular}{cccc}
\toprule
Step & Smallest time remaining & Placement rule & Partial sequence \\
\midrule
1 & $2$ on $J_4$ for $M_1$ & place first & $J_4,\_,\_,\_,\_$ \\
2 & $3$ on $J_2$ for $M_1$ & next from front & $J_4,J_2,\_,\_,\_$ \\
3 & $4$ on $J_3$ for $M_2$ & place last & $J_4,J_2,\_,\_,J_3$ \\
4 & $5$ on $J_1$ for $M_2$ & next from back & $J_4,J_2,\_,J_1,J_3$ \\
5 & $5$ on $J_5$ for $M_1$ & fill remaining front slot & $J_4,J_2,J_5,J_1,J_3$ \\
\bottomrule
\end{tabular}
\end{center}

\item Gantt timings:

\begin{center}
\begin{tabular}{lcc}
\toprule
Job & Machine 1 & Machine 2 \\
\midrule
$J_4$ & 0--2 & 2--11 \\
$J_2$ & 2--5 & 11--19 \\
$J_5$ & 5--10 & 19--25 \\
$J_1$ & 10--17 & 25--30 \\
$J_3$ & 17--23 & 30--34 \\
\bottomrule
\end{tabular}
\end{center}

Minimum makespan:
\[
34
\]

\item The original order $J_1,J_2,J_3,J_4,J_5$ gives makespan 39.

So Johnson's Rule improves the schedule by
\[
39-34=5
\]
time units.

\item Solving the job shop problem gives an optimal schedule with makespan
\[
19
\]

One optimal schedule is:

\begin{center}
\begin{tabular}{lcccc}
\toprule
Job & Task & Machine & Start & End \\
\midrule
$A$ & 1 & $M_1$ & 0 & 4 \\
$A$ & 2 & $M_3$ & 5 & 10 \\
$A$ & 3 & $M_2$ & 13 & 16 \\
$B$ & 1 & $M_2$ & 0 & 3 \\
$B$ & 2 & $M_1$ & 6 & 10 \\
$B$ & 3 & $M_3$ & 10 & 16 \\
$C$ & 1 & $M_3$ & 0 & 5 \\
$C$ & 2 & $M_2$ & 5 & 7 \\
$C$ & 3 & $M_1$ & 10 & 14 \\
$D$ & 1 & $M_1$ & 4 & 6 \\
$D$ & 2 & $M_2$ & 7 & 13 \\
$D$ & 3 & $M_3$ & 16 & 19 \\
\bottomrule
\end{tabular}
\end{center}

This can be obtained using a job shop formulation or an OR-Tools interval-variable model with precedence constraints within each job and \texttt{NoOverlap} constraints on each machine.

\item Add a due-date constraint for job $C$:
\[
C_C \le 15
\]
where $C_C$ is the completion time of job $C$.

In the schedule above, job $C$ already finishes at time 14, so this additional constraint is satisfied and one optimal makespan-19 schedule remains feasible. In other instances, such a due-date constraint could force resequencing or even cause infeasibility.
\end{enumerate}

\textbf{2. Binary Integer Programming}

\begin{enumerate}
\item Let $x_i=1$ if base $i$ is opened and 0 otherwise.
Let $y_j=1$ if district $j$ is covered and 0 otherwise.

\[
\max Z = 9000y_1+12000y_2+14000y_3+10000y_4+11000y_5+8000y_6
\]
subject to
\[
14x_A+11x_B+13x_C+10x_D+12x_E+9x_F \le 31
\]
\[
x_A+x_B+x_C+x_D+x_E+x_F \le 3
\]
\[
y_1 \le x_A+x_B
\]
\[
y_2 \le x_A+x_C+x_E
\]
\[
y_3 \le x_B+x_C+x_F
\]
\[
y_4 \le x_A+x_D+x_F
\]
\[
y_5 \le x_C+x_D+x_E
\]
\[
y_6 \le x_E+x_F
\]
\[
x_i,y_j \in \{0,1\}
\]

\item The optimal solution is
\[
x_A=x_C=1,\qquad x_B=x_D=x_E=x_F=0
\]
Total cost:
\[
14+13=27
\]
Covered districts:
\[
1,2,3,4,5
\]
Total covered population:
\[
9000+12000+14000+10000+11000=56000
\]

\item Yes. In the optimal solution above, district 6 remains uncovered. This is acceptable in a maximum-coverage model because the objective is to maximise covered population subject to tight resources. In a set covering model, complete coverage would be mandatory.

\item Revised model changes:

District 6 must be covered:
\[
y_6 = 1
\]
or equivalently
\[
x_E+x_F \ge 1
\]

Bases $B$ and $F$ cannot both be chosen:
\[
x_B+x_F \le 1
\]

Under these revised conditions, one optimal solution is
\[
x_E=x_F=1,\qquad x_A=x_B=x_C=x_D=0
\]
with covered districts $2,3,4,5,6$ and total covered population
\[
12000+14000+10000+11000+8000=55000
\]
There are multiple optimal solutions with the same covered population; for example, $x_C=x_F=1$ is also optimal.
\end{enumerate}

\textbf{3. Simulation}

\begin{enumerate}
\item Main components:

Entities: students.

Resource: advisers.

Events: arrival, service start, service completion.

Stochastic inputs: interarrival and service times.

Measures: number served, average wait, maximum wait, percentage waiting less than 8 minutes.

\item One run with seed 24 gave:

\begin{center}
\begin{tabular}{lc}
\toprule
Measure & Result \\
\midrule
Students served & 41 \\
Average wait & 1.17 minutes \\
Maximum wait & 10.32 minutes \\
\% waiting less than 8 minutes & 95.12\% \\
\bottomrule
\end{tabular}
\end{center}

\item Peak-period results:

\begin{center}
\begin{tabular}{lcccc}
\toprule
Advisers & Served & Avg wait & Max wait & \% below 8 min \\
\midrule
2 & 70 & 1.88 & 12.99 & 95.71\% \\
3 & 73 & 0.32 & 3.22 & 100\% \\
4 & 70 & 0.08 & 2.23 & 100\% \\
\bottomrule
\end{tabular}
\end{center}

From this run, the minimum staffing level meeting the target is 3 advisers.

\item A \texttt{PriorityResource} can be used in place of a standard \texttt{Resource}. Each arriving student would be labelled priority or standard, for example with probability 0.10 of being priority, and then would request service with a smaller priority number if they are priority users.

In SimPy terms, the main change is from a standard request such as
\[
\texttt{with desk.request() as req:}
\]
to a priority request such as
\[
\texttt{with desk.request(priority=0) as req:}
\]
for priority users and a larger number such as \texttt{priority=1} for standard users.

This would usually reduce priority-user waits and increase the waits experienced by standard users.

\item Replications are needed because a single run may be unusually good or bad. If the model represented a continuing system rather than a finite session, a warm-up period might also be appropriate.
\end{enumerate}

\textbf{4. Forecasting}

\begin{enumerate}
\item Using periods 1--12 only, centred 4-quarter moving averages are available for periods 3--10:

\begin{center}
\begin{tabular}{ccc}
\toprule
Period & Demand & CMA \\
\midrule
3 & 470 & 491.250 \\
4 & 560 & 499.500 \\
5 & 450 & 507.750 \\
6 & 536 & 516.500 \\
7 & 500 & 525.875 \\
8 & 600 & 535.125 \\
9 & 485 & 544.375 \\
10 & 575 & 554.375 \\
\bottomrule
\end{tabular}
\end{center}

A linear trend fitted to these centred moving averages is
\[
\hat{T}_t = 463.21 + 9.02t
\]
where $t=1$ corresponds to 2019 Q1.

\item Seasonal indices:

Additive:
\[
Q1=-58.72,\quad Q2=19.91,\quad Q3=-23.72,\quad Q4=62.53
\]

Multiplicative:
\[
Q1=0.8884,\quad Q2=1.0372,\quad Q3=0.9535,\quad Q4=1.1209
\]

\item Validation forecasts for 2022:

\begin{center}
\begin{tabular}{lccc}
\toprule
Quarter & Actual & Additive Forecast & Multiplicative Forecast \\
\midrule
2022 Q1 & 522 & 521.76 & 515.68 \\
2022 Q2 & 618 & 609.41 & 611.44 \\
2022 Q3 & 572 & 574.80 & 570.71 \\
2022 Q4 & 694 & 670.07 & 680.99 \\
\bottomrule
\end{tabular}
\end{center}

Using MAD:
\[
\text{MAD}_{\text{add}} \approx 8.89,\qquad \text{MAD}_{\text{mult}} \approx 6.80
\]
where absolute error is calculated as $|\text{Actual} - \text{Forecast}|$.

For the multiplicative model:
\[
|522-515.68|=6.32,\quad |618-611.44|=6.56,\quad |572-570.71|=1.29,\quad |694-680.99|=13.01
\]
so
\[
\text{MAD}_{\text{mult}}=\frac{6.32+6.56+1.29+13.01}{4}\approx 6.80
\]

For the additive model:
\[
|522-521.76|=0.24,\quad |618-609.41|=8.59,\quad |572-574.80|=2.80,\quad |694-670.07|=23.93
\]
so
\[
\text{MAD}_{\text{add}}=\frac{0.24+8.59+2.80+23.93}{4}\approx 8.89
\]
The multiplicative model is better.

\item Forecasts for 2023 using the multiplicative model:

\begin{center}
\begin{tabular}{lc}
\toprule
Quarter & Forecast \\
\midrule
2023 Q1 & 547.74 \\
2023 Q2 & 648.86 \\
2023 Q3 & 605.11 \\
2023 Q4 & 721.43 \\
\bottomrule
\end{tabular}
\end{center}
\end{enumerate}

\textbf{5. Linear Programming}

Let $C$ be casings and $I$ be inserts.

\begin{enumerate}
\item A suitable model is
\[
\max Z = 40C + 34I
\]
subject to
\[
4C + 2I \le 100
\]
\[
C + 3I \le 90
\]
\[
I \le 24
\]
\[
C,I \ge 0
\]

\item The optimal solution is
\[
C=13,\qquad I=24
\]
Maximum contribution:
\[
Z=40(13)+34(24)=1336
\]

\item Binding constraints: moulding and insert demand cap.

Slack values:
\[
4(13)+2(24)=100 \Rightarrow 0 \text{ moulding slack}
\]
\[
13+3(24)=85 \Rightarrow 5 \text{ assembly hours unused}
\]
\[
24-24=0
\]

\item The new condition is
\[
I \ge 0.25(C+I)
\]
which simplifies to
\[
C-3I \le 0
\]
At the original optimum,
\[
13-3(24)=-59 \le 0
\]
so the original solution remains feasible and optimal.

\item With 12 extra assembly hours, the assembly constraint becomes less restrictive, but the current solution already has 5 hours of assembly slack and is limited by moulding plus the insert cap. The optimal plan therefore does not change.

Extra contribution gained:
\[
0
\]
Maximum cost of buying all 12 hours:
\[
12 \times 3 = 36
\]
So the extra hours should not be bought.
\end{enumerate}

\textbf{6. Transportation}

\begin{enumerate}
\item Let $x_{ij}$ be the quantity shipped from warehouse $i$ to store $j$.

A suitable model is
\[
\min Z = \sum_i \sum_j c_{ij}x_{ij}
\]
subject to the supply constraints for each warehouse $i$:
\[
\sum_j x_{ij} = s_i
\]
demand constraints for each store $j$:
\[
\sum_i x_{ij} = d_j
\]
and non-negativity:
\[
x_{ij}\ge 0
\]

\item One optimal plan is:

\begin{center}
\begin{tabular}{lc}
\toprule
Route & Quantity \\
\midrule
$W_1 \to S_1$ & 20 \\
$W_1 \to S_4$ & 60 \\
$W_2 \to S_2$ & 60 \\
$W_2 \to S_3$ & 40 \\
$W_3 \to S_1$ & 50 \\
$W_3 \to S_3$ & 40 \\
\bottomrule
\end{tabular}
\end{center}

Total transportation cost:
\[
1560
\]

\item With $W_2 \not\to S_2$ and at least 20 units from $W_1$ to $S_4$, one optimal revised plan is:

\begin{center}
\begin{tabular}{lc}
\toprule
Route & Quantity \\
\midrule
$W_1 \to S_2$ & 60 \\
$W_1 \to S_4$ & 20 \\
$W_2 \to S_3$ & 60 \\
$W_2 \to S_4$ & 40 \\
$W_3 \to S_1$ & 70 \\
$W_3 \to S_3$ & 20 \\
\bottomrule
\end{tabular}
\end{center}

Revised cost:
\[
1680
\]
Increase in cost:
\[
1680-1560=120
\]

\item Example extension: add vehicle capacities or delivery time windows so the model reflects route feasibility as well as shipping cost.
\end{enumerate}

\end{document}
